\documentclass[11pt]{article}
\usepackage[margin=1in]{geometry}
\usepackage{graphicx}
\usepackage{booktabs}
\usepackage{amsmath,amssymb}
\usepackage{enumitem}
\usepackage{hyperref}
\usepackage[table]{xcolor}
\usepackage{caption}
\usepackage{subcaption}

\definecolor{grokgreen}{HTML}{2E8B57}
\definecolor{nogrokred}{HTML}{CC3333}

\hypersetup{colorlinks=true, linkcolor=blue, citecolor=blue, urlcolor=blue}

\title{Experiment 6: Mechanistic Analysis of Grokking under Federated Learning}
\author{Federated Learning Grokking Project}
\date{March 2026}

\begin{document}
\maketitle

% =============================================================================
\section{Overview}
\label{sec:exp6-overview}

\begin{itemize}
  \item \textbf{Goal:} Post-hoc mechanistic analysis of 237 FL runs from
        Exps~2--3 to answer: \emph{why} does grokking survive (or fail)
        under federated averaging?
  \item \textbf{No new training runs.} All analysis uses per-round metrics
        already logged: test/train accuracy, loss, weight norms (layers~1
        and~2), IPR, mean client drift $\bar{d}_k$, and client weight
        divergence $\sigma_w$.
  \item \textbf{Dataset:} 191 grokked runs + 46 failed runs across
        $\alpha \in \{0.2, 0.25, 0.3, 0.35, 0.5\}$, $K \in \{2, 5, 10,
        20, 50, 97\}$, and $\alpha_{\mathrm{dir}} \in \{0.01\text{--}1000\}$.
\end{itemize}

\noindent\textbf{Four main findings:}
\begin{enumerate}[label=(\roman*)]
  \item Drift predicts grokking delay \emph{within} each $\alpha$, not
        merely as a proxy for $\alpha$.
  \item IPR onset consistently precedes drift drop (72/72 runs, 12/12
        threshold combinations).
  \item Accuracy, IPR, and normalized drift collapse onto shared curves
        under rescaled time.
  \item FL disrupts generalization but not memorization: $T_{\mathrm{memo}}$
        is invariant to $K$ while $T_{\mathrm{grok}}$ varies by up to
        $2.1\times$.
\end{enumerate}


% =============================================================================
\section{Representative Run Comparison (Figure~\ref{fig:representative})}
\label{sec:representative}

\noindent Four runs spanning the mechanistic spectrum:

\begin{itemize}
  \item \textbf{Healthy grok:} $\alpha{=}0.5$, IID, $K{=}10$
        --- fast grokking at ${\sim}7{,}600$ steps
  \item \textbf{Delayed grok:} $\alpha{=}0.25$, $\alpha_{\mathrm{dir}}{=}0.01$,
        $K{=}10$ --- groks at ${\sim}36{,}600$ steps
  \item \textbf{Failed (data-limited):} $\alpha{=}0.2$, IID, $K{=}10$
        --- fails centralized too
  \item \textbf{Failed (fragmented):} $\alpha{=}0.25$, IID, $K{=}97$
        --- too many clients
\end{itemize}

\noindent\textbf{Key observations from six-panel comparison:}

\begin{itemize}
  \item \textbf{Failed runs} show persistently high drift ($>0.8$),
        flat IPR (${\sim}0.03$), and train loss that plateaus.
        Weight norms never reach grokking-associated magnitudes.
  \item \textbf{Delayed grok} follows the same trajectory as the healthy
        grok (IPR rises, drift drops, weights plateau) but stretched in time.
        Final values are identical.
  \item \textbf{Mean client drift} (panel~d) separates all four cases most
        cleanly: $K{=}97$ failure $>0.8$; healthy grok $<0.3$.
\end{itemize}

\begin{figure}[htbp]
  \centering
  \includegraphics[width=\textwidth]{figures/fig6_representative_comparison.png}
  \caption{\textbf{Representative run comparison.}
  (a)~Test accuracy, (b)~train loss (log), (c)~IPR,
  (d)~mean client drift, (e)~client weight divergence,
  (f)~layer-1 weight norm.}
  \label{fig:representative}
\end{figure}


% =============================================================================
\section{Drift Predicts Grokking Delay, Controlling for $\alpha$
  (Figure~\ref{fig:drift-controlled})}
\label{sec:drift}

\subsection{The confound}

\begin{itemize}
  \item Naive analysis: $\rho = 0.70$, $p = 1.6 \times 10^{-29}$
        across all 237 runs.
  \item \textbf{But:} 42 of 46 failed runs are $\alpha{=}0.2$, which fails
        in the centralized setting too.
  \item Drift correlates with $\alpha$ because both are controlled by the
        same experimental knobs.
  \item \textbf{Honest question:} does drift predict $T_{\mathrm{grok}}$
        \emph{within a fixed $\alpha$}?
\end{itemize}

\subsection{Within-$\alpha$ correlation (panel~b)}

\begin{itemize}
  \item Within-$\alpha$ Spearman rank correlations between drift and
        $T_{\mathrm{grok}}$:
\end{itemize}

\begin{center}
\begin{tabular}{lcccc}
\toprule
$\alpha$ & $n$ & Spearman $\rho$ & $p$-value & Significance \\
\midrule
0.25 & 47 & 0.79 & $< 10^{-3}$ & $^{***}$ \\
0.30 & 51 & 0.70 & $< 10^{-3}$ & $^{***}$ \\
0.35 & 42 & 0.52 & $< 10^{-3}$ & $^{***}$ \\
0.50 & 51 & 0.34 & $< 0.05$    & $^{*}$   \\
\bottomrule
\end{tabular}
\end{center}

\begin{itemize}
  \item Strong near the boundary ($\rho = 0.79$ at $\alpha{=}0.25$),
        weakens in the interior ($\rho = 0.34$ at $\alpha{=}0.5$).
  \item Expected if drift modulates grokking primarily near criticality.
\end{itemize}

\subsection{Causal chain: $K \to$ drift $\to T_{\mathrm{grok}}$ (panel~c)}

\begin{itemize}
  \item At $\alpha{=}0.3$ (IID, Exp~2): increasing $K$ from 2 to 97
        monotonically increases both drift and $T_{\mathrm{grok}}$.
  \item Heterogeneity is controlled (IID), so the only variable is $K$.
  \item $\rho = 1.0$ between drift and $T_{\mathrm{grok}}$ across $K$ values.
\end{itemize}

\begin{figure}[htbp]
  \centering
  \includegraphics[width=\textwidth]{figures/fig7_drift_alpha_controlled.png}
  \caption{\textbf{Drift predicts grokking delay, controlling for $\alpha$.}
  (a)~$T_{\mathrm{grok}}$ vs.\ drift (all runs; red box: 42/46 failures
  are $\alpha{=}0.2$).
  (b)~Within-$\alpha$ Spearman $\rho$.
  (c)~$K \to$ drift $\to T_{\mathrm{grok}}$ at $\alpha{=}0.3$.}
  \label{fig:drift-controlled}
\end{figure}


% =============================================================================
\section{Temporal Ordering: IPR Rise Precedes Drift Drop
  (Figure~\ref{fig:temporal})}
\label{sec:temporal}

\subsection{Question}

\begin{itemize}
  \item Do clients align first, then Fourier features emerge?
        (drift drops $\to$ IPR rises)
  \item Or do Fourier features crystallize first, then clients converge?
        (IPR rises $\to$ drift drops)
\end{itemize}

\subsection{Method}

\begin{itemize}
  \item For each of 72 grokking runs (Exp~3a, $K{=}10$), detect:
  \begin{itemize}
    \item $T_{\mathrm{IPR\ onset}}$: first round IPR exceeds $2.5\times$
          its early-training baseline
    \item $T_{\mathrm{drift\ drop}}$: first round (post-peak) smoothed
          drift falls below 50\% of its peak
  \end{itemize}
\end{itemize}

\subsection{Results}

\begin{itemize}
  \item IPR onset precedes drift drop in \textbf{72 of 72} runs
        (zero exceptions).
  \item Median lag: $T_{\mathrm{drift\ drop}} - T_{\mathrm{IPR\ onset}}
        = 10{,}775$ steps.
  \item Wilcoxon signed-rank test: $p = 1.8 \times 10^{-13}$.
\end{itemize}

\subsection{Robustness}

\begin{itemize}
  \item Tested 12 combinations: IPR threshold
        $\in \{2\times, 2.5\times, 3\times, 4\times\}$ $\times$
        drift fraction $\in \{30\%, 50\%, 70\%\}$.
  \item \textbf{All 12 settings:} IPR rises first in 100\% of runs
        (Figure~\ref{fig:temporal}d).
\end{itemize}

\subsection{Caveats}

\begin{itemize}
  \item \textbf{Temporal precedence $\neq$ causation.} Both events could
        be driven by a third variable (e.g., weight norm growth reaching
        a critical value).
  \item \textbf{Threshold asymmetry.} IPR onset = crossing a low absolute
        value (${\sim}0.05$); drift drop = falling below a fraction of peak.
        The \emph{ordering} is robust; the \emph{lag magnitude} depends on
        threshold choice.
\end{itemize}

\begin{figure}[htbp]
  \centering
  \includegraphics[width=\textwidth]{figures/fig8_temporal_ordering.png}
  \caption{\textbf{Temporal ordering.}
  (a)~$T_{\mathrm{IPR\ onset}}$ vs.\ $T_{\mathrm{drift\ drop}}$:
  all points above diagonal.
  (b)~Lag histogram (median = 10{,}775 steps).
  (c)~Box plots: both precede $T_{\mathrm{grok}}$; IPR leads more.
  (d)~Robustness heatmap: 100\% IPR-first across all threshold combos.}
  \label{fig:temporal}
\end{figure}


% =============================================================================
\section{Scaling Collapse (Figure~\ref{fig:collapse})}
\label{sec:collapse}

\subsection{Setup}

\begin{itemize}
  \item Rescale time as $\tau = (t - T_{\mathrm{grok}}) /
        (T_{\mathrm{grok}} - T_{50})$, where $T_{50}$ = first step
        test accuracy reaches 50\%.
  \item Apply to 117 grokking runs from Exps~2--3.
\end{itemize}

\subsection{Accuracy collapse (panel~b) --- expected}

\begin{itemize}
  \item Test accuracy curves collapse onto a sigmoid (fitted $k \approx 2.48$).
  \item \textbf{Partly tautological:} centering on the transition and scaling
        by the width will collapse any sigmoid-like curve.
  \item Confirms transitions are sigmoidal, but not evidence of deeper
        universality on its own.
\end{itemize}

\subsection{IPR and drift collapse (panels~c, d) --- non-trivial}

\begin{itemize}
  \item IPR \textbf{also collapses} under the same rescaling, despite not
        being constrained by the procedure:
  \begin{itemize}
    \item Begins rising at $\tau \approx -3$ (3 transition widths before
          grokking)
    \item Plateaus at $\tau \approx 0$
  \end{itemize}
  \item Normalized drift follows a universal inverted-U:
  \begin{itemize}
    \item Peaks at $\tau \approx -2$
    \item Declines through the transition
    \item Stabilizes at ${\sim}50\%$ of peak after grokking ($\tau > 0$)
  \end{itemize}
  \item \textbf{Implication:} feature crystallization, client convergence,
        and generalization evolve on a single intrinsic timescale that varies
        across configurations but maintains fixed internal structure.
\end{itemize}

\begin{figure}[htbp]
  \centering
  \includegraphics[width=\textwidth]{figures/fig9_universal_collapse.png}
  \caption{\textbf{Scaling collapse.}
  (a)~Raw trajectories (117 runs).
  (b)~Rescaled accuracy (expected for sigmoids).
  (c)~IPR collapse (non-trivial): features emerge at $\tau \approx -3$.
  (d)~Drift collapse (non-trivial): peaks at $\tau \approx -2$.
  Black = median; shading = IQR.}
  \label{fig:collapse}
\end{figure}


% =============================================================================
\section{Two-Timescale Separation (Figure~\ref{fig:two-timescales})}
\label{sec:timescales}

\subsection{Definitions}

\begin{itemize}
  \item $T_{\mathrm{memo}}$: first step with train accuracy $\geq 99\%$
        (memorization complete)
  \item $T_{\mathrm{grok}}$: first step with test accuracy $\geq 95\%$
        permanently (generalization)
  \item Generalization delay: $\Delta T = T_{\mathrm{grok}} - T_{\mathrm{memo}}$
\end{itemize}

\subsection{Within-$\alpha$ variability (panel~a)}

\begin{itemize}
  \item CV of $T_{\mathrm{memo}}$ vs.\ $T_{\mathrm{grok}}$ within each
        $\alpha$, across all $K$ and heterogeneity:
\end{itemize}

\begin{center}
\begin{tabular}{lccc}
\toprule
$\alpha$ & CV($T_{\mathrm{memo}}$) & CV($T_{\mathrm{grok}}$) & Ratio \\
\midrule
0.25 & 0.042 & 0.159 & $4\times$ \\
0.30 & 0.033 & 0.068 & $2\times$ \\
0.35 & 0.022 & 0.026 & $1\times$ \\
0.50 & 0.005 & 0.005 & $1\times$ \\
\bottomrule
\end{tabular}
\end{center}

\begin{itemize}
  \item Near the boundary ($\alpha{=}0.25$): $T_{\mathrm{grok}}$ is
        $4\times$ more variable than $T_{\mathrm{memo}}$.
  \item In the interior ($\alpha \geq 0.35$): both equally stable.
\end{itemize}

\subsection{$T_{\mathrm{memo}}$ vs.\ $K$ (panel~b)}

\begin{itemize}
  \item $T_{\mathrm{memo}}$ is nearly flat from $K{=}2$ to $K{=}50$:
  \begin{itemize}
    \item $\alpha{=}0.3$: within $\pm 1\%$
    \item $\alpha{=}0.5$: within $\pm 0.3\%$
    \item Only at $K{=}97$: modest $+17\%$ increase ($\alpha{=}0.3$)
  \end{itemize}
  \item $T_{\mathrm{grok}}$ diverges:
  \begin{itemize}
    \item $\alpha{=}0.3$: $12{,}822$ ($K{=}2$) $\to$ $27{,}422$ ($K{=}97$)
          --- $2.1\times$ slowdown
  \end{itemize}
  \item \textbf{Conclusion:} FL does not impair memorization. It specifically
        disrupts the transition from memorization to generalization.
\end{itemize}

\subsection{Drift prolongs generalization delay (panel~c)}

\begin{itemize}
  \item $\Delta T = T_{\mathrm{grok}} - T_{\mathrm{memo}}$ vs.\ drift,
        with within-$\alpha$ Spearman:
  \begin{itemize}
    \item $\alpha{=}0.25$: $\rho \approx 0.73$
    \item Weakens in the interior, paralleling Section~\ref{sec:drift}.
  \end{itemize}
\end{itemize}

\begin{figure}[htbp]
  \centering
  \includegraphics[width=\textwidth]{figures/fig10_two_timescales.png}
  \caption{\textbf{FL disrupts generalization, not memorization.}
  (a)~Within-$\alpha$ CV: $T_{\mathrm{grok}}$ up to $4\times$ more
  variable than $T_{\mathrm{memo}}$.
  (b)~$T_{\mathrm{memo}}$ (dashed, flat) vs.\ $T_{\mathrm{grok}}$
  (solid, diverging) vs.\ $K$.
  (c)~Generalization delay vs.\ drift with within-$\alpha$ correlations.}
  \label{fig:two-timescales}
\end{figure}


% =============================================================================
\section{Summary}
\label{sec:exp6-summary}

\begin{enumerate}
  \item \textbf{Drift predicts grokking delay within each $\alpha$.}
  \begin{itemize}
    \item Within-$\alpha$ Spearman $\rho$: 0.34 ($\alpha{=}0.5$) to
          0.79 ($\alpha{=}0.25$)
    \item Not a proxy for $\alpha$; carries independent predictive information
    \item \emph{Caveat:} 42/46 failures are $\alpha{=}0.2$ (fails centralized)
  \end{itemize}

  \item \textbf{IPR rise precedes drift drop (temporal ordering).}
  \begin{itemize}
    \item 72/72 runs, 12/12 threshold combinations, zero exceptions
    \item Consistent with: Fourier features crystallize in global model first,
          then clients converge
    \item \emph{Caveat:} temporal ordering $\neq$ causation; lag magnitude
          depends on threshold choice
  \end{itemize}

  \item \textbf{Approximate scaling collapse.}
  \begin{itemize}
    \item Under $\tau = (t - T_{\mathrm{grok}})/(T_{\mathrm{grok}} - T_{50})$:
          test accuracy, IPR, and drift collapse onto shared curves
    \item Accuracy collapse is partly tautological (any sigmoid collapses)
    \item IPR and drift collapses are non-trivial $\to$ single intrinsic
          timescale governs the transition
  \end{itemize}

  \item \textbf{FL disrupts generalization, not memorization.}
  \begin{itemize}
    \item $T_{\mathrm{memo}}$ invariant to $K$ (CV $\approx$ 0.005--0.04)
    \item $T_{\mathrm{grok}}$ varies by up to $2.1\times$
    \item Drift specifically prolongs the memorization-to-generalization gap
  \end{itemize}
\end{enumerate}

\subsection{Limitations}

\begin{itemize}
  \item \textbf{No model checkpoints.} Cannot compute Fourier spectra of
        $W_1$ at specific timepoints to test whether clients converge to
        the same Fourier modes. Requires re-running with checkpoint saving.
  \item \textbf{No per-client trajectories.} Individual client IPR, loss,
        and weights are aggregated before logging. Cannot distinguish
        synchronized vs.\ leader-follower transitions.
  \item \textbf{Limited FL-specific failures.} Only 4/46 failures are
        FL-caused (not $\alpha{=}0.2$). Exps~4--5 will provide more.
  \item \textbf{Single architecture and task.} All results are for 2-layer
        MLP on modular addition ($p{=}97$). Generalization to other settings
        is open.
\end{itemize}

\end{document}
